\section*{Chapter \ref{ch:b2:oceanic-islands} --- Oceanic Islands}

\begin{solution}{exo:b2:oceanic-islands:1}
(a) Preshield stage, like Loihi today. (b) \omterm{def:b2:oceanic-islands:stages}{Shield stage}. (c) Postshield stage. (d) Rejuvenated
stage. (e) No eruptive stage any more: the volcano has sunk under its reef, which is now an \omterm{def:b2:oceanic-islands:reefs}{atoll}.
\end{solution}

\begin{solution}{exo:b2:oceanic-islands:2}
(a) $2500 + 350\sqrt9 = \qty{3550}{m}$. (b) $350\times(\sqrt{25} - \sqrt9) = 350\times2 =
\qty{700}{m}$. (c) $350/(2\times5) = \qty{35}{m}$ per million years.
\end{solution}

\begin{solution}{exo:b2:oceanic-islands:3}
(a) $5000/17 = \qty{294}{kyr}$ around the main islands; $30\,000/70 = \qty{429}{kyr}$ along the
ridge. (b) Slides are missed: a younger volcano can bury the scar and the debris of an older one,
and old deposits are covered by sediment. The true number is larger and the true interval
shorter.
\end{solution}

\begin{solution}{exo:b2:oceanic-islands:4}
(a) $4158\times0.3048 = \qty{1267}{m}$; $4610\times0.3048 = \qty{1405}{m}$; $150\times0.3048 =
\qty{46}{m}$. (b) The first reef grew no deeper than \qty{46}{m} and its base is now at
\qty{1267}{m}: the volcano has sunk by at least $1267 - 46 = \qty{1221}{m}$ (more if sea level
was higher then). (c) The sinking is at most about \qty{1267}{m}, over at least \qty{33.9}{Myr}:
at most $1267/33.9 = \qty{37}{m}$ per million years, \qty{0.037}{mm} a year.
\end{solution}

\begin{solution}{exo:b2:oceanic-islands:5}
(a) $\dot S = 175/\sqrt t$: \qty{87.5}{m.Myr^{-1}} (\qty{0.0875}{mm.yr^{-1}}) at \qty{4}{Myr},
\qty{35}{m.Myr^{-1}} (\qty{0.035}{mm.yr^{-1}}) at \qty{25}{Myr}, \qty{25}{m.Myr^{-1}}
(\qty{0.025}{mm.yr^{-1}}) at \qty{49}{Myr}. (b) $4.2/0.0875 = 48$ times. (c) $2.6/0.0875 = 30$
times: the loading phase sinks an island much faster than the cooling of a young plate.
\end{solution}

\begin{solution}{exo:b2:oceanic-islands:6}
(a) $\dot S = (3200/62.8)\,e^{-100/62.8} = 50.96\times0.2035 = \qty{10.37}{m}$ per million years,
\qty{0.0104}{m} per thousand years: the measured \qty{0.01}{m} per thousand years. (b) $e^{-t/62.8}
= 10/50.96$, $t = 62.8\ln5.096 = \qty{102.3}{Myr}$. (c) Root law: $2500 + 350\times5 =
\qty{4250}{m}$; exponential law: $6400 - 3200\times0.6716 = \qty{4251}{m}$; they differ by
\qty{1}{m}.
\end{solution}

\begin{solution}{exo:b2:oceanic-islands:7}
(a) $\qty{8.6}{cm.yr^{-1}} = \qty{86}{km.Myr^{-1}}$; $86\times33 = \qty{2838}{km}$. (b) Yes:
\qty{27.7}{Myr} is less than the \qty{33}{Myr} after which the reefs drown, so its reef should
still be alive. (c) $55\times1.609 = \qty{88.5}{km}$, or $88.5/86 = \qty{1.0}{Myr}$ more: about
\qty{28.7}{Myr}. That is somewhat less than \qty{33}{Myr}, but the \qty{33}{Myr} is itself an
approximate figure for the whole chain, and growth also depends on the water temperature at each
\omterm{def:b2:oceanic-islands:reefs}{atoll}: Kure, the last living \omterm{def:b2:oceanic-islands:reefs}{atoll}, is indeed close to the \omterm{def:b2:oceanic-islands:darwin-point}{Darwin point}.
\end{solution}

\begin{solution}{exo:b2:oceanic-islands:8}
(a) Box: $30\times17\times2 = \qty{1020}{km^3}$; pyramid: a third, \qty{340}{km^3}. (b) One block
alone is a tenth (box) or a thirtieth (pyramid) of the whole volcano of Kilauea. (c) $23\,000\times
0.1 = \qty{2300}{km^3}$, a quarter of Kilauea; an estimate that stands or falls with the assumed
thickness.
\end{solution}

\begin{solution}{exo:b2:oceanic-islands:9}
(a) $1000/2.6 = \qty{385}{kyr}$. (b) Yes: \qty{2.6}{m} per thousand years is
\qty{2.6}{mm.yr^{-1}}, less than \qty{4.2}{mm.yr^{-1}}, a margin of \qty{1.6}{mm.yr^{-1}}.
(c) Relative rise $12 + 2.6 = \qty{14.6}{mm.yr^{-1}}$; the reef falls behind by $14.6 - 4.2 =
\qty{10.4}{mm.yr^{-1}}$ for $16.5 - 8.2 = \qty{8.3}{kyr}$: its top ends at $10.4\times8.3 =
\qty{86}{m}$. (d) $(12 - 4.2)\times8.3 = \qty{65}{m}$. In both cases it stays above the deepest
growth limit but in much dimmer water: whether it catches up depends on how its growth slows with
depth.
\end{solution}

\begin{solution}{exo:b2:oceanic-islands:10}
(a) $S(t_1) = C(\sqrt{t_1 + \Delta} - \sqrt{t_1})$ has the derivative $\frac C2(1/\sqrt{t_1 +
\Delta} - 1/\sqrt{t_1}) < 0$: it decreases with $t_1$. (b) $350\sqrt{20} = \qty{1565}{m}$; $350\times
(\sqrt{70} - \sqrt{50}) = 350\times1.2955 = \qty{453}{m}$. (c) No: $\sqrt t$ is finite at $0$, so
the sinking over any interval is finite, \qty{1565}{m} in (b). Only the rate at the instant $t = 0$
is infinite, where the law is no longer meant to describe the floor in detail.
\end{solution}

\begin{solution}{exo:b2:oceanic-islands:11}
(a) $30\times0.1 = \qty{3}{m}$. (b) About \qty{120}{m} if its top stood at sea level at the
highstand: the \omterm{def:b2:oceanic-islands:reefs}{atoll} becomes a limestone island, rain dissolves the aragonite of the corals and the
rock recrystallises to calcite, with caves. (c) $120/3 = 40$ cycles, \qty{4}{Myr}. (d) The leached
zones are these times of exposure, sealed under the reef that grew again when the sea came back.
Over many cycles the sea returns to similar highstands, so $\Delta\eta$ is small, the reef top
regrows to sea level at each highstand, and the thickness of the cap is set by the sinking, as the
proposition says.
\end{solution}

\begin{solution}{exo:b2:oceanic-islands:12}
(a) $3200\times(e^{-50/62.8} - e^{-120/62.8}) = 3200\times(0.4510 - 0.1480) = \qty{970}{m}$.
(b) $150 + 970 = \qty{1120}{m}$. (c) $700\times1.8288 = \qty{1280}{m}$, \qty{160}{m} deeper than the
estimate. (d) The \omterm{def:b2:oceanic-islands:seamount}{guyot} may have drowned earlier, on a younger plate, and sunk more; a neighbouring
volcano may have loaded it; its top may have been cut lower by waves, or be covered by sediment; the
plate under it may be of a different age; and the sea level of the time of drowning was not
today's.
\end{solution}

\begin{solution}{pb:b2:oceanic-islands:1}
\textbf{1.} $2500 + 350\sqrt{10} = 2500 + 350\times3.1623 = \qty{3607}{m}$.
\textbf{2.} $3607 + 1000 = \qty{4607}{m}$.
\textbf{3.} $V = \pi r^2h/3 = \pi\times50^2\times4.607/3 = \qty{12061}{km^3}$.
\textbf{4.} $0.85\times12\,061 = \qty{10252}{km^3}$ in \qty{1}{Myr}: \qty{0.0103}{km^3} a year,
ten times less than Kilauea today.
\textbf{5.} $95\,600/12\,061 = 7.9$: about eight.
\textbf{6.} Tholeiitic lava, from the \omterm{def:b2:oceanic-islands:stages}{shield stage}; the alkalic lavas of the postshield stage cap
it.
\textbf{7.} $\qty{1}{km}/\qty{1}{Myr} = \qty{1}{m}$ per thousand years, 2.6 times slower than off
the Island of Hawaii.
\textbf{8.} Yes: $4.2 - 1.0 = \qty{3.2}{mm.yr^{-1}}$ to spare.
\textbf{9.} $12 + 1 - 4.2 = \qty{8.8}{mm}$ a year.
\textbf{10.} $150/8.8 = \qty{17}{kyr}$. A drowned reef terrace on the flank, a step of the kind
found off the Island of Hawaii.
\textbf{11.} $350\times(\sqrt{22} - \sqrt{10}) = 350\times(4.6904 - 3.1623) = \qty{535}{m}$.
\textbf{12.} $1000 - 535 = \qty{465}{m}$, at $465/12 = \qty{38.8}{m}$ per million years.
\textbf{13.} $350/(2\sqrt{10}) = \qty{55.3}{m.Myr^{-1}}$ at first; $350/(2\sqrt{22}) =
\qty{37.3}{m.Myr^{-1}}$ at the end.
\textbf{14.} $2500 + 350\times4.6904 = \qty{4142}{m}$.
\textbf{15.} The large Hawaiian islands rise up to about four times higher, so there is far more
rock above sea level to remove; their \omterm{def:b2:oceanic-islands:stages}{shield stage} also lasted longer.
\textbf{16.} Waves, which cut back the coast and plane a wave-cut platform; landslides and \omterm{def:b2:oceanic-islands:collapse}{flank
collapses}.
\textbf{17.} A \omterm{def:b2:oceanic-islands:reefs}{barrier reef}, with a lagoon between the reef and the last of the volcanic shore.
\textbf{18.} $350\times(\sqrt{70} - \sqrt{22}) = 350\times(8.3666 - 4.6904) = \qty{1287}{m}$.
\textbf{19.} $1287/48 = \qty{26.8}{m}$ per million years, \qty{0.027}{mm} a year: $4.2/0.0268 =
157$ times less than the reef could grow.
\textbf{20.} \qty{20}{m} thicker than the limestone over the basalt under Parry Island
(\qty{1267}{m}), and \qty{118}{m} thinner than under Elugelab (\qty{1405}{m}): the same size.
\textbf{21.} $2500 + 350\times8.3666 = \qty{5428}{m}$.
\textbf{22.} $3200\times(e^{-70/62.8} - e^{-120/62.8}) = 3200\times(0.3280 - 0.1480) =
\qty{576}{m}$.
\textbf{23.} $150 + 576 = \qty{726}{m}$.
\textbf{24.} The reef cap of Petrel \omterm{def:b2:oceanic-islands:reefs}{atoll} is \textbf{\qty{1287}{m}} thick, about \qty{1.3}{km}
of reef limestone, as under Enewetak.
\end{solution}
